 %!TEX root = main.tex
 
\paragraph{Most damaging scenario.}
Scenario~S3, the simultaneous energy crisis and Singapore hub
closure,  is the most destructive disruption in the scenario
library, with $\Delta Z_F = +\$3{,}989{,}134$ ($+89.6\%$ above
$Z^*$). Scenario~S1 (hub closure alone) is a close second at
$+81.6\%$. Both are an order of magnitude more destructive than
the worst purely price-based disruption: T3 (hub surcharge)
costs $+19.1\%$, and even the Suez corridor crisis~S4 reaches
only $+7.7\%$. The ratio of worst topological to worst price
disruption is $89.6\,/\,19.1 \approx 4.7\times$. The asymmetry
is structural: a price shock raises the cost of existing routes
uniformly and the network re-optimises around them; a topological
shock disconnects suppliers from the network entirely, leaving
no routing alternative and forcing spot-market procurement at
a $2.5\times$ premium. Geographic chokepoints (Suez Canal,
$+7.7\%$) are manageable; topological chokepoints (H3,
$+81.6\%$) are existential.
 
\paragraph{Persistent structural weakness.}
Two structural vulnerabilities recur across the entire analysis.
The first and most severe is the exclusive dependence of
four suppliers on hub~H3~(Singapore): S4~(Taipei),
S5~(Seoul), S6~(Yokohama), and S9~(Chengdu) have no alternative
hub connection, making H3 a hard single point of failure for
48.7\% of the supply base. The min-cut analysis
(Section~\ref{sec:phase1}) quantifies this: removing H3 collapses
network throughput by 46\%, more than three times the impact of
any other node. The second vulnerability is subtler: the same
ten inter-warehouse arcs (A185, A188, A190, A199, A210, A223,
A224, A241, A251, A317) are independently re-discovered by
strategy~A as natural relief paths under T1, T3, S1, and~S3.
Their systematic recurrence is a signal that the Phase~1
optimisation horizon is too narrow: it minimises operating cost
under normal conditions while leaving an adaptation tax of up to
\$98{,}244 (T3) payable each time a shock forces reconfiguration.
 
\paragraph{Cost of resilience.}
Section~\ref{sec:resilient_baseline} provides a direct, quantified
answer. Forcing open the ten insurance arcs and adding four
synthetic bypass arcs, one per stranded supplier, priced at
$\times 1.5$ above a conservative distance estimate,
raises $Z^*$ by \$10{,}513 ($+0.24\%$). Under a Singapore hub
closure (S1), this \$10{,}513 investment reduces the disruption
cost from $+81.6\%$ to $+41.3\%$ and saves \$1{,}776{,}606,
delivering a resilience ratio of \$169 saved per \$1 invested.
A network designed with robustness as a secondary objective
alongside cost minimisation would, at negligible premium,
eliminate the primary structural vulnerability identified
throughout this analysis. The cost of building resilience is
two orders of magnitude smaller than the cost of experiencing
the disruption it protects against.